\documentclass{article}

\usepackage[nonatbib]{neurips_2026}

\usepackage[utf8]{inputenc}
\usepackage[T1]{fontenc}
\usepackage{hyperref}
\usepackage{url}
\usepackage{booktabs}
\usepackage{amsfonts}
\usepackage{nicefrac}
\usepackage{microtype}
\usepackage{xcolor}
\usepackage{graphicx}
\usepackage{amsmath}

\title{SMC-QA: Selective Memory Consolidation with Query-Aware Retrieval for Long-Context Question Answering}

\author{
  Zheng Yifan \\
  SID: 12532245 \\
  \texttt{yifan.zheng@example.edu} \\
  \And
  Chen Zeming \\
  SID: 12531410 \\
  \texttt{zeming.chen@example.edu} \\
}

\begin{document}

\maketitle

\begin{abstract}
Long-context question answering requires systems to manage and retrieve information from extensive interaction histories. A flat memory approach treats all past information equally, which may introduce noise and miss the distinction between recent and consolidated knowledge. In this project, we propose SMC-QA, a lightweight hierarchical memory framework that organizes information into short-term memory (STM) and long-term memory (LTM). We design a selective promotion strategy based on relevance, reuse frequency, and information diversity to consolidate valuable short-term memories into long-term storage, along with rule-based query-aware retrieval routing. We evaluate SMC-QA and five baselines on the LongMemEval benchmark using both oracle and full-haystack settings. Our results show that (1) hierarchical memory consistently outperforms single-layer memory, (2) selective promotion reduces redundancy compared to naive merging, and (3) the diversity component of our promotion score contributes the most to retrieval quality among the three scoring factors.
\end{abstract}

%======================================================================
\section{Introduction}
%======================================================================

Long-context question answering (QA) is a challenging task that requires systems to locate relevant evidence within large volumes of historical interactions. Unlike standard QA, where the context is typically short and self-contained, long-context QA involves multi-session dialogues, knowledge updates, and temporal reasoning across extended interaction histories.

A fundamental challenge in this setting is \textbf{memory management}: how should a system decide what information to retain, what to discard, and where to search when a new question arrives? A simple approach is to store all past information in a single pool and retrieve from it directly, but this leads to increased noise as irrelevant content accumulates.

Inspired by cognitive science concepts of short-term and long-term memory, as well as recent work on hierarchical memory systems for LLMs~\cite{himem, rmm, lightmem}, we propose \textbf{SMC-QA} (Selective Memory Consolidation for Question Answering). Our system:

\begin{itemize}
    \item Organizes information into a \textbf{short-term memory} (STM) for recent content and a \textbf{long-term memory} (LTM) for consolidated knowledge.
    \item Uses a lightweight \textbf{promotion score} combining relevance, reuse frequency, and diversity to decide which STM items should be promoted to LTM.
    \item Applies \textbf{rule-based query-aware routing} to decide whether to search STM, LTM, or both based on query characteristics.
\end{itemize}

We compare SMC-QA against five baselines on the LongMemEval benchmark~\cite{longmemeval} and conduct ablation studies to understand the contribution of each component.

%======================================================================
\section{Related Work}
%======================================================================

\textbf{LongMemEval}~\cite{longmemeval} provides a benchmark for evaluating long-term interactive memory across five capabilities: information extraction, multi-session reasoning, knowledge updates, temporal reasoning, and abstention. We use this benchmark as our primary evaluation dataset.

\textbf{HiMem}~\cite{himem} introduces hierarchical long-term memory for long-horizon dialogue agents, supporting both hierarchical memory construction and hybrid retrieval. Our STM/LTM structure draws inspiration from this work but uses a simpler, non-parametric design.

\textbf{Reflective Memory Management (RMM)}~\cite{rmm} proposes multi-granularity memory summarization across utterances, turns, and sessions. While we do not implement summary-based promotion in our current system, this work motivates our consideration of memory consolidation.

\textbf{LightMem}~\cite{lightmem} modularizes memory into retrieval, writing, and long-term consolidation components. Our system adopts a similar modular philosophy but focuses specifically on the promotion decision and retrieval routing.

%======================================================================
\section{Method}
%======================================================================

\subsection{System Overview}

SMC-QA processes a stream of text chunks from historical interactions and organizes them into a two-level memory hierarchy. When a query arrives, the system routes it to the appropriate memory level(s) and retrieves the most relevant evidence.

\subsection{Memory Structure}

\textbf{Short-Term Memory (STM)} operates as a fixed-capacity FIFO buffer (capacity $C_s = 40$ chunks). New chunks enter STM, and when capacity is exceeded, the oldest chunks are evicted. STM preserves recent, detailed information.

\textbf{Long-Term Memory (LTM)} stores promoted items with a maximum capacity of $C_l = 120$ chunks. Items enter LTM only through the promotion process. When LTM exceeds capacity, the least-accessed items are evicted. A duplicate detection mechanism (cosine similarity $\geq 0.85$) prevents redundant entries.

\subsection{Selective Memory Consolidation}

Every $f = 5$ chunks, a consolidation step evaluates STM items for promotion to LTM. Each candidate receives a composite score:

\begin{equation}
    \text{score} = w_r \cdot \text{relevance} + w_u \cdot \text{reuse} + w_d \cdot \text{diversity}
\end{equation}

where $w_r = 0.4$, $w_u = 0.4$, $w_d = 0.2$.

\textbf{Relevance} measures the average cosine similarity between the chunk embedding and the embeddings of recent queries (window of 5).

\textbf{Reuse} captures how frequently the chunk has been retrieved, normalized by the maximum hit count in STM.

\textbf{Diversity} measures novelty relative to existing LTM content: $\text{diversity} = 1 - \max_{l \in \text{LTM}} \text{sim}(c, l)$.

Only the top-$k$ candidates ($k=3$) with scores above a threshold ($\tau = 0.35$) are promoted.

\subsection{Query-Aware Retrieval Routing}

When a query arrives, the system applies three rules in sequence:

\begin{enumerate}
    \item \textbf{Rule 1 (LTM-first)}: If the query contains temporal keywords (e.g., ``before'', ``earlier'', ``previously'', ``first''), route to LTM.
    \item \textbf{Rule 2 (STM-first)}: If the query's average similarity to recent STM chunks exceeds its similarity to LTM by a margin $\delta = 0.05$, route to STM.
    \item \textbf{Rule 3 (Hybrid)}: Otherwise, retrieve from both STM and LTM, merge results, and re-rank.
\end{enumerate}

For single-path retrieval, we return top-4 results. For hybrid retrieval, we take top-4 from each level, merge, deduplicate, and re-rank to return top-6.

%======================================================================
\section{Experimental Setup}
%======================================================================

\subsection{Dataset}

We use the \textbf{LongMemEval} benchmark~\cite{longmemeval}, which contains 500 evaluation instances across six question types: single-session-user, single-session-assistant, single-session-preference, multi-session, temporal-reasoning, and knowledge-update.

We evaluate in two settings:
\begin{itemize}
    \item \textbf{Oracle}: Each instance includes only answer-relevant sessions ($\sim$28 turns/instance).
    \item \textbf{S-Cleaned (Full Haystack)}: Each instance includes $\sim$50 sessions with $\sim$500 turns, where answer-relevant content is mixed with distractors. We truncate to 200 turns per instance for computational tractability.
\end{itemize}

For oracle experiments, we sample 200 instances (50 dev / 150 test) with seed 42 and run formal experiments with three seeds (42, 43, 44).

\subsection{Baselines}

\begin{enumerate}
    \item \textbf{Flat Memory}: All chunks in one pool; retrieve top-6 by cosine similarity.
    \item \textbf{STM-only}: Retrieve only from the most recent 40 chunks.
    \item \textbf{LTM-only}: Build LTM via our promotion method; retrieve only from LTM.
    \item \textbf{Naive STM+LTM}: Build STM/LTM; always retrieve from both (no routing).
    \item \textbf{Frequency-based Promotion}: Promote to LTM based on retrieval frequency only.
\end{enumerate}

\subsection{Metrics}

\begin{itemize}
    \item \textbf{Hit@6}: Whether at least one gold evidence turn is found in the top-6 retrieved chunks.
    \item \textbf{Recall@6}: Fraction of gold evidence turns covered by the top-6 retrieved chunks.
    \item \textbf{Contains}: Whether the concatenated top-3 retrieved chunks contain the gold answer string.
\end{itemize}

We match retrieved chunks to gold evidence using word-level overlap with a threshold of 50\%.

\subsection{Implementation}

We use \texttt{all-MiniLM-L6-v2} from sentence-transformers for embedding (384-dimensional). All retrieval uses brute-force cosine similarity. Text is chunked at 120 words with 30-word overlap, preserving natural turn boundaries when turns are shorter than 160 words.

%======================================================================
\section{Results}
%======================================================================

\subsection{Oracle Setting}

Table~\ref{tab:main} shows results on the oracle dataset (150 test instances, mean $\pm$ std across 3 seeds).

\begin{table}[h]
\centering
\caption{Main results on LongMemEval Oracle (150 test instances, 3 seeds).}
\label{tab:main}
\begin{tabular}{lccc}
\toprule
\textbf{Method} & \textbf{Hit@6} & \textbf{Recall@6} & \textbf{Contains} \\
\midrule
STM-only & 0.642 $\pm$ 0.003 & 0.478 $\pm$ 0.002 & 0.342 $\pm$ 0.042 \\
LTM-only & 0.713 $\pm$ 0.019 & 0.549 $\pm$ 0.014 & 0.309 $\pm$ 0.044 \\
Flat Memory & 0.782 $\pm$ 0.019 & 0.623 $\pm$ 0.021 & 0.298 $\pm$ 0.049 \\
Naive STM+LTM & 0.776 $\pm$ 0.017 & 0.614 $\pm$ 0.018 & 0.329 $\pm$ 0.039 \\
Freq Promotion & 0.804 $\pm$ 0.013 & 0.641 $\pm$ 0.021 & 0.327 $\pm$ 0.045 \\
SMC-QA (Ours) & 0.778 $\pm$ 0.022 & 0.612 $\pm$ 0.018 & \textbf{0.333 $\pm$ 0.030} \\
\bottomrule
\end{tabular}
\end{table}

\textbf{Key observations:}
\begin{itemize}
    \item \textbf{Hierarchical memory helps}: STM-only (0.642 Hit@6) significantly underperforms all methods that include LTM, confirming that relying solely on recent memory is insufficient.
    \item \textbf{LTM adds value}: LTM-only (0.713) outperforms STM-only by 7.1\%, showing that promoted content provides useful signal even without short-term context.
    \item \textbf{Promotion strategies are competitive}: Both Freq Promotion (0.804) and SMC-QA (0.778) perform comparably to or better than Flat Memory (0.782) on Hit@6, despite having access to only a subset of chunks through their memory structures.
    \item \textbf{Contains metric}: SMC-QA achieves the highest Contains score (0.333) with the lowest variance (0.030), suggesting it retrieves more answer-relevant content.
\end{itemize}

\subsection{Full Haystack Setting (S-Cleaned)}

Table~\ref{tab:scleaned} shows results when answer-relevant content is mixed with $\sim$200 distractor turns.

\begin{table}[h]
\centering
\caption{Results on LongMemEval S-Cleaned (100 instances, 200 turns each).}
\label{tab:scleaned}
\begin{tabular}{lccc}
\toprule
\textbf{Method} & \textbf{Hit@6} & \textbf{Recall@6} & \textbf{Contains} \\
\midrule
STM-only & 0.050 & 0.033 & 0.070 \\
LTM-only & 0.420 & 0.290 & 0.270 \\
Flat Memory & \textbf{0.460} & \textbf{0.315} & 0.260 \\
Naive STM+LTM & 0.420 & 0.290 & 0.250 \\
Freq Promotion & 0.430 & 0.292 & 0.250 \\
SMC-QA (Ours) & 0.420 & 0.290 & 0.250 \\
\bottomrule
\end{tabular}
\end{table}

In the full haystack setting, \textbf{STM-only collapses} to near-zero performance (0.050 Hit@6), demonstrating that fixed-window recent memory is completely inadequate when answer evidence may lie far in the past. Flat Memory achieves the highest scores here because it can search across all chunks without the information loss introduced by STM's capacity limit. However, the gap between hierarchical methods and Flat Memory is small (4\% Hit@6), and hierarchical methods have the advantage of bounded memory size.

\subsection{Ablation Study}

Table~\ref{tab:ablation} shows the effect of removing each component from SMC-QA.

\begin{table}[h]
\centering
\caption{Ablation results on LongMemEval Oracle (150 test instances).}
\label{tab:ablation}
\begin{tabular}{lccc}
\toprule
\textbf{Configuration} & \textbf{Hit@6} & \textbf{Recall@6} & \textbf{Contains} \\
\midrule
SMC-QA (full) & 0.753 & 0.588 & 0.327 \\
$-$ relevance & 0.760 & 0.589 & 0.333 \\
$-$ diversity & 0.740 & 0.570 & 0.313 \\
$-$ routing & 0.753 & 0.589 & 0.320 \\
\bottomrule
\end{tabular}
\end{table}

\textbf{Diversity matters most}: Removing the diversity component causes the largest drop ($-$1.3\% Hit@6, $-$1.8\% Recall@6), indicating that preventing duplicate information from entering LTM is the most valuable aspect of our promotion strategy.

\textbf{Relevance has minimal impact}: Removing relevance slightly improves performance, suggesting that in the oracle setting where most content is already relevant, this signal provides less discriminative power.

\textbf{Routing has limited effect}: Removing query-aware routing shows no change in Hit@6 and slight decrease in Contains, suggesting that the rule-based routing provides modest benefits in this dataset.

\subsection{Case Analysis}

We analyze specific instances where SMC-QA outperforms or underperforms Flat Memory.

\textbf{Success cases} tend to involve \textbf{knowledge-update} questions (e.g., ``Where did I go on my most recent family trip?''). In these cases, SMC-QA's LTM retains the updated information while Flat Memory may rank outdated mentions higher due to more repetitions.

\textbf{Failure cases} tend to involve \textbf{specific factual details} (e.g., ``How many plants did I initially plant?''). Here, the relevant information may not have been promoted to LTM because it appeared infrequently and had low reuse scores.

%======================================================================
\section{Discussion}
%======================================================================

\textbf{When does hierarchical memory help?} Our results show that the benefit of hierarchical memory is most clear when (1) the conversation history is long enough that STM alone cannot cover relevant information, and (2) the question specifically requires long-term knowledge. In the oracle setting with limited context, the advantage is modest.

\textbf{Promotion strategy trade-offs}: The frequency-based baseline performs surprisingly well, suggesting that simple heuristics can be effective. Our multi-factor promotion score (relevance + reuse + diversity) shows its strength primarily through the diversity component, which prevents LTM from accumulating redundant entries.

\textbf{Limitations}: (1) Our routing rules are hand-crafted and may not generalize to all query types. A learned router could be more adaptive. (2) We do not perform summary-based promotion, which could compress information more effectively. (3) The embedding model (MiniLM-L6-v2) is relatively small; stronger encoders may improve retrieval quality. (4) We evaluate primarily on retrieval metrics; integrating an answer generation model would provide a more complete picture.

\textbf{Future work}: Extending this framework with a learned routing module, summary-based consolidation, and answer generation evaluation would address these limitations and provide stronger evidence for the effectiveness of hierarchical memory in long-context QA.

%======================================================================
\section{Conclusion}
%======================================================================

We presented SMC-QA, a lightweight hierarchical memory framework for long-context question answering. Through systematic comparison with five baselines on the LongMemEval benchmark, we showed that: (1) hierarchical memory (STM+LTM) consistently outperforms single-layer memory; (2) selective promotion strategies (both frequency-based and multi-factor) improve over naive memory merging; and (3) the diversity component of our promotion score is the most impactful factor. While Flat Memory remains competitive in our evaluation settings, hierarchical approaches offer the practical advantage of bounded memory size and become increasingly important as conversation histories grow longer.

%======================================================================
\section{Contributions}
%======================================================================

\textbf{Zheng Yifan} (50\%): System design, implementation of memory structures, promotion strategies, routing module, experiment pipeline, and report writing.

\textbf{Chen Zeming} (50\%): Data processing, baseline implementation, evaluation metrics, ablation experiments, and case analysis.

%======================================================================
% References
%======================================================================

\begin{thebibliography}{4}

\bibitem{longmemeval}
Wu, D., Wang, H., Yu, W., Zhang, Y., Chang, K.-W., and Yu, D.
\newblock LongMemEval: Benchmarking Chat Assistants on Long-Term Interactive Memory.
\newblock \emph{arXiv preprint arXiv:2410.10813}, 2024.

\bibitem{himem}
Zhang, N., Yang, X., Tan, Z., Deng, W., and Wang, W.
\newblock HiMem: Hierarchical Long-Term Memory for LLM Long-Horizon Agents.
\newblock \emph{arXiv preprint arXiv:2601.06377}, 2026.

\bibitem{rmm}
Tan, Z., Yan, J., Hsu, I.-H., Han, R., Wang, Z., Le, L., Song, Y., Chen, Y., Palangi, H., Lee, G., Iyer, A.~R., Chen, T., Liu, H., Lee, C.-Y., and Pfister, T.
\newblock In Prospect and Retrospect: Reflective Memory Management for Long-term Personalized Dialogue Agents.
\newblock In \emph{Proceedings of ACL}, pages 8416--8439, 2025.

\bibitem{lightmem}
Fang, J., Deng, X., Xu, H., Jiang, Z., Tang, Y., Xu, Z., Deng, S., Yao, Y., Wang, M., Qiao, S., Chen, H., and Zhang, N.
\newblock LightMem: Lightweight and Efficient Memory-Augmented Generation.
\newblock \emph{arXiv preprint arXiv:2510.18866}, 2025.

\end{thebibliography}

\end{document}
